% Independent analytic audit, 2026-10-09. No numerical certificate is rerun.
% Fragment only; all pre-existing manuscripts and certificates are unchanged.
\subsection{A noncircular compiler for the causal ordinary lower start}
\label{sec:own-lower-compiler-chain-audit-20261009}

This audit concerns the ordinary starting terms of the project's own
Chen lower bound. It does not certify the double-sieve action bank, the
eleven-term pointwise inequality, or a new numerical constant. In
particular it does not obtain a lower bound from the positive-product
upper constant \(K\). The original sequence here is
\[
 {\cal A}_N=\{N-p:p<N,\ p\ {\rm prime}\},\qquad N\ {\rm even}.
\]
We use
\[
 a=\frac{417}{5000},\quad h=\frac{19}{125},\quad
 \nu=\frac{249}{448},\quad \Lambda=\frac{153}{256},\quad
 \mu=\frac{25}{128},\quad \sigma=\frac7{32}.
\]
The conclusions below apply separately to \(w=a\) and \(w=h\), and
indeed to every fixed \(0<w\le h\). They apply to the finite causal
policies in \texttt{lower-start-joint-finite-approximation.tex}, with
their complete branch labels and margins, not to arbitrary bounded
arithmetic restrictions or arbitrary five-dimensional choices of caps.

\paragraph{1. What needed to be supplied.}
The approximation lemma explicitly assumes a finite-policy signed
compiler. The starting-lower note describes the sign reversal, whereas
\texttt{actual-k-compilation-joint-audit.tex} proves the upper version.
Also the older end of \texttt{prefix-sieve-interface.tex} states its
ordinary density discussion only on compact \(s\ge2\). That sentence
alone is insufficient at the negative fifth node, whose upper residual
parameter can lie in \(1<s<2\).

There are genuine proofs, rather than numerical PASS records, behind
these steps: \texttt{retained-base-upper.tex} proves the lower residual
equality; \texttt{fourth-prefix-K-refinement-route.tex} supplies the
upper counterpart; the two actual-\(K\) joint audits prove atom
compilation and distribution separately. We give the lower compilation
and the actual sequence mapping below. We also strengthen the density
argument by an exact prime-deletion comparison. This closes the above
range and assembly omissions without assuming a two-sided dimension
hypothesis uniform in the target \(N\).

\paragraph{2. The actual finite lower tree.}
Fix first all finite policy data and a small root buffer
\(\rho>0\); put \(\nu_\rho=\nu-\rho\), with \(2h<\nu_\rho\).
On decreasing primes \(p_i=N^{t_i}\), start with the ordinary lower
support at level \(N^{\nu_\rho}\). Its initial test is \(p_1^2<N^{\nu_\rho}\),
which never fails below \(N^w\); its other tests are
\(S_{j-1}+3t_j<\nu_\rho\) at even \(j\).
The finite state records all ancestor decisions, the selected branch,
and its actual frozen cap.

An admitted third prefix obeys \(t_1+3t_2<\nu_\rho\).
Freeze \(\lambda_3\ge\nu_\rho\) there, changing subsequent even tests.
The fourth prefix is admitted only if
\[
 S_3+3t_4<\lambda_3.                                      \tag{A}
\]
Only after this test may it freeze \(\lambda_4\ge\lambda_3\).
Every \(0<t_5<t_4\) is then admitted, and a fifth state may freeze
\(\lambda_5\ge\lambda_4\). Subsequent even tests use \(\lambda_5\).
An inherited label means literal equality of caps and no new test.
Restoring the inherited rule on an unused or boundary cell is allowed;
deleting that node is not an equivalent operation.

Thus every first exclusion has even full rank. For any finite prime
set \(Q\) and any \(0\le q(p)\le1\), the exact first-exit identity is
\begin{equation}\label{eq:own-lower-even-exit}
 \sum_{d\mid Q}\mu(d){\bf1}_{d\in{\cal T}}q(d)
 =\prod_{p\mid Q}(1-q(p))
 -\sum_{P=p_1\cdots p_j\ {\rm first\ exit}}
       q(P)\prod_{\substack{p\mid Q\\p<p_j}}(1-q(p)).
\end{equation}
Partition the excluded divisors by their unique first excluded prefix;
its sign is \(+1\), and subtract its complete continuation. This proves
the identity, also for inherited history-dependent levels.
With \(q(p)={\bf1}_{p\mid n}\), it proves a pointwise lower weight.
For comparison, retaining \(1\) but excluding a single prime gives
the value \(1>0\) at that prime: an extra odd first exclusion would
invalidate this argument. The present construction makes none.

\paragraph{3. Full and selected support, including inherited states.}
Every complete descendant is in the full square support at its
current cap. Initially
\[
 2t_1<\nu_\rho,\qquad t_1+2t_2<t_1+3t_2<\nu_\rho,\qquad
 S_2+2t_3<t_1+3t_2<\nu_\rho.
\]
At rank four its own square follows from (A), and at rank five
\begin{equation}\label{eq:own-lower-fifth-square}
 S_4+2t_5<S_3+3t_4<\lambda_3\le\lambda_4\le\lambda_5.
\end{equation}
Thereafter an even cube implies its own square and the following odd
square, using
\(S_{j-1}+2t_j\le S_{j-2}+3t_{j-1}\).
All products have exponent below \(\Lambda+o(1)<3/5\).

For a selected nonempty subset \(I\subset\{1,\ldots,r\}\) put
\[
 b_I=\sum_{i\in I}t_i,\qquad
 G_I=\max_{\substack{j\le r\\j\notin I}}
 \left(b_I+\sum_{\substack{i<j\\i\notin I}}t_i+2t_j\right).
\]
An active selected branch must freeze its cap with strict slack above
the parent, every full square and \(G_I\), and strictly below
\[
 v_\alpha(b_I)=
 \min\left\{\frac{1+b_I}{2},
 \frac{1-(3/2-2\sigma)b_I}{1+\sigma},
 \frac{2-\alpha\sigma-3b_I}{2}\right\}.
\]
These are whole-box tests, not just a test on \(b_I\).
For every later slot the selected-remaining square is exactly the full
square \(S_{j-1}+2t_j\), because every selected slot has already occurred.
An inherited branch therefore continues to license every descendant.
At rank three a \(W\) branch keeps its first-three-slot ordering and
all its original gates; later primes are smaller. The two-capacity
and \(W\)-branch repairs in the actual compilation audit apply to
square-supported slot vectors, independently of the sign of their
coefficients. In particular its \(C<1\) reassignment of all slots is
needed; mere existence of a selected product is not sufficient.

The untouched ordinary tree is covered too. For \(t_1<h<\mu\),
the repaired full envelope satisfies
\[
 \Phi(t_1,t_2,t_3)\ge
 \max\left\{\frac{1+t_1}{2},\Lambda-\frac{3t_1}{8}\right\}
 \ge\nu.
\]
The affine expressions intersect at \(25/224\). If a raw \(W\) gate
is inactive, its single-variable rescue is retained in this maximum.
A source factorization parameter strictly between \(\nu_\rho\) and
\(\nu\) supplies the ordinary branch margin.

In all branches factorize the upper endpoint vector of each sufficiently
fine labelled box at each Type II scale. Keep that partition of the
slots fixed on the whole box; normalize each group's convolution by
its number-of-slots factorial. Coordinatewise reduction preserves
the required monomials (notably \(5-2\sigma>0\)). This is the
coefficient-level factorization required by dispersion. It is stronger
than factoring each individual integer by a different rule.

\paragraph{4. The lower compiler, with both inequalities displayed.}
Write \(g_N(p)={\bf1}_{p>2,\ p\nmid N}/(p-1)\) and
\(V_N(z)=\prod_{p<z}(1-g_N(p))\). Choose a presieve threshold
\(z_0=N^\omega\) below every active prefix cutoff. Let
\(R_{\cal T}\) be the raw lower density restricted to rough divisors
with primes in \([z_0,N^w)\), and \(V_{\cal T}\) its full raw density.
Equation \eqref{eq:own-lower-even-exit} gives
\begin{equation}\label{eq:own-lower-rough-projection}
 V_N(z_0)R_{\cal T}-V_{\cal T}
 =\sum_{\substack{P\ {\rm full\ first\ exit}\\p_{\rm last}<z_0}}
       g_N(P)V_N(p_{\rm last})\ge0.
\end{equation}
Every rough first exit is the same full first exit with its last prime
at least \(z_0\); all its preceding primes are then rough as well.

Refine fixed coarse labels into disjoint half-open prime intervals.
Use inner tests for positive even-rank coefficients and outer tests
for negative odd-rank coefficients, retaining every ancestor test.
The resulting rough coefficients satisfy \(c(d)\le\lambda_{\cal T}(d)\).
Fine-box discrepancies occupy finitely many cube/gate strips; at
fixed \(\omega\) their weighted mass tends to zero with the fine mesh.
The tagged-prime estimate of the compilation audit is unchanged in
absolute value:
\[
 V_N(z_0)\sum_d|c(d)-\lambda_{\cal T}(d)|g_N(d)
 \ll V_N(N^w)C_{\cal T}\tau\omega^{-5}+o_N(V_N(N^w)).
\]
Coarse discontinuities are first buffered with the inherited rule.
This is not interpolation across a cap jump.

Let \(B^-\le I_0\le B^+\) be the small-prime beta polynomials, where
\(I_0(n)={\bf1}_{(n,P_N(z_0))=1}\).
Attach \(B^-\) when \(c(d)>0\) and \(B^+\) when \(c(d)<0\).
For each divisor separately its contribution is at most \(c(d)I_0(n)\).
Consequently the compiled coefficient \(\rho_N\) satisfies
\begin{equation}\label{eq:own-lower-compiled-pointwise}
 \sum_{q\mid n}\rho_N(q)
 \le I_0(n)\sum_{\substack{d\mid n\\d\ {\rm rough}}}c(d)
 \le{\bf1}_{(n,P_N(N^w))=1}.
\end{equation}
No nonnegativity of \(B^-\) is assumed. The unique small/rough
factorization gives \(|\rho_N(q)|\le1\). The beta density error is at
most \(O(\varepsilon_\beta\omega^{-2}V_N(N^w))\). Combining this with
\eqref{eq:own-lower-rough-projection}, rather than reversing it, gives
\[
 \sum_q\rho_N(q)g_N(q)\ge
 V_{\cal T}-O((C_{\cal T}\tau\omega^{-5}
              +\varepsilon_\beta\omega^{-2}+o_N(1))V_N(N^w)).
\]

For a retained rough label vector \(J\), with repeated-label
multiplicities \(m_I\), its exact atom is
\[
 \frac{(-1)^{|J|}}{\prod_I m_I!}\,
 {\bf1}_{\rm squarefree}
   (\chi_{J_1}*\cdots*\chi_{J_{|J|}})*\beta_J.
\]
The squarefree indicator applies to the rough product before beta.
The scalar is retained once. Partitioning its slots into three groups
and dividing by the three group factorials gives one-bounded group
coefficients; multiplying those factorials into the scalar restores
the same atom. Attach beta only to the third group. Its small primes
cannot collide with rough primes. For fixed policy, cutoff and mesh
the atom count and total absolute scalar mass are finite independently
of \(N\). The standard beta fundamental lemma is the external
construction input here; the sign assembly above is a project proof.

\paragraph{5. Mapping these atoms to the original prime sequence.}
This is the \(k=1\) case of the signed distribution derivation in
\texttt{actual-k-distribution-joint-audit.tex}, not an application
of its positive \(K\) conclusion. On a local prime scale
\(x\ge N^{1-\kappa}/O(1)\), the residue is exactly \(N\).
Choose all local losses smaller than the fixed branch margin and
\(\kappa<\epsilon/(2(1+\epsilon))\), so \(N<x^{1+\epsilon}\)
eventually. One may use the buffered residue exponent
\(\alpha=1+3\kappa\) in the factorization gates.

The primary inputs, read in the saved L23 v1 text, are Definition 2.1,
Proposition 5.2, Proposition 3.2 and Lemma 3.4; Proposition 6.1 supplies
the original branch factorization. They are statements in
\href{https://arxiv.org/abs/2309.08522v1}{Lichtman's 2023 preprint}.
Their spectral/dispersive proofs are external inputs, not numerical
facts and not relabelled here as a checked journal publication.
Proposition 5.2 allows signed complex one-bounded coefficients and
a supremum over \(0<|a|<x^{1+\epsilon}\).

For completeness the actual matching inequalities are as follows.
The order-three Heath--Brown expansion of the one original prime
variable has at most six slots. A slot with length
\(x^{2\zeta}\le U\le x^{1/3+\zeta/2}\) is used as the SW side.
Put \(A_0=Ux^{-\zeta}\). The fixed slot factorization gives
\[
 B_1\le A_0,\quad A_0^2B_2B_3^2\le x^{1-3\zeta},\quad
 N^\sigma A_0B_1B_2^5B_3^2\le x^{2-3\zeta},\quad
 A_0^{1+2\sigma}B_1B_2^{5-2\sigma}B_3^{2+2\sigma}
       \le x^{2-2\zeta}.
\]
Replacing \(A_0\) by \(U\) leaves savings respectively
\(\zeta,2\zeta,(1-2\sigma)\zeta=9\zeta/16\).
After the still smaller beta and box losses, these imply exactly
\[
 Q_1<Ux^{-\epsilon},\quad U^2Q_2Q_3^2<x^{1-8\epsilon},\quad
 N^\sigma UQ_1Q_2^5Q_3^2<x^{2-15\epsilon},\quad
 U^{1+2\sigma}Q_1Q_2^{5-2\sigma}Q_3^{2+2\sigma}<x^{2-15\epsilon}.
\]
The individual rough \(1\) or rough \(\mu\) slot has the auxiliary
coprimality SW property, including interval restrictions and the
smooth twists, by the project's explicit rough-SW proof. The
HB rough threshold \(x^{1/(\log\log x)^3}\) here is distinct from
the modulus presieve threshold \(z_0=N^\omega\).
The
divisor-bounded complementary coefficient is truncated and normalized
as in the distribution audit; its discarded part is controlled by
shifted-divisor moments. An arbitrary arithmetic restriction on the
SW slot would not be allowed.

If no medium slot exists, preservation of the original group scale
forces at least one long slot; there are at most two. A long slot
cannot be a truncated Mobius slot and hence has coefficient \(1\)
before partial summation. One long slot is covered by Lemma 3.4:
its length is at least \(x^{1-10\zeta}\), exceeding the whole modulus
times \(x^\epsilon\). Two long slots are covered by Proposition 3.2:
the short product is at most \(x^{20\zeta}\), so the modulus exponent
\(<3/5\) is below \((1-20\zeta)(2/3-3\epsilon)\).
The extension to separate arbitrary long-slot intervals is explicitly
in the proof of Proposition 5.2 and Lemma 7.1 of
\href{https://arxiv.org/abs/2006.07088v1}{Maynard II v1};
it retains the supremum over primitive residues.

Rough ranks at most two here have total support
\(N^{2h+O(\ell+\tau)}\), with \(2h=38/125<1/2\); classical
Bombieri--Vinogradov handles them after choosing \(\kappa,\ell,\tau\)
small. This assertion is specific to the present ordinary-start
atoms; it is not the low-rank check for a larger marked-prime
double-sieve family.

Removing rough-square collisions costs
\((N+N^{\theta_*})N^{-\omega+o(1)}\), \(\theta_*<3/5\).
This is applied to the actual coefficients, or with the fixed divisor
bound after HB; an unweighted \(N/d+1\) bound is not used alone on
unbounded HB coefficients. The short-prime tail
\(p<2N^{1-\kappa}\) costs \(N^{1-\kappa+o(1)}\) before distribution,
since the divisor sum is bounded by \(\tau(N-p)\). Its exact-coprime
mean has only a logarithmic additional factor. Thus a near-\(N\)
local-scale theorem really covers the complete prime sequence.
Dyadic decompositions and Mellin tails cost fixed logarithmic powers.

It follows, for this same fixed \(\rho_N\), that
\begin{equation}\label{eq:own-lower-signed-prime-remainder}
 \sum_{\substack{q\\(q,N)=1}}\rho_N(q)
 \left(\pi(N;q,N)-\frac{\pi(N)}{\varphi(q)}\right)
 =O_A(N\log^{-A}N).
\end{equation}
The source's exact coprime mean differs from \(\pi(N)/\varphi(q)\)
by at most \(\omega(q)/\varphi(q)\) for a prime sequence; summed over
\(q\le N^{\theta_*}\) this is \(O(\log^2N)\).
Every modulus prime is restricted to be coprime to \(N\).
All constants may depend on the fixed finite policy. No assertion
of distribution for a policy or a positive cutoff varying with \(N\)
is required.

\paragraph{6. Exact ordinary residual densities, uniformly in \(N\).}
Define ordinary upper support by the cube tests at odd ranks, and
ordinary lower support by \(q_1^2<R\) and the cube tests at even
ranks. Write their actual multiplicative density sums as \(V_N^\pm(R,y)\).
For the lower support assume \(y\le\sqrt R\), so its first square
does not create an odd exit. We prove, on compact positive parameter
sets \(R=N^\beta,y=N^t\),
\begin{equation}\label{eq:own-lower-ordinary-equalities}
 \frac{V_N^+(N^\beta,N^t)}{V_N(N^t)}=F(\beta/t)+o(1)
       \quad(\beta/t>1),\qquad
 \frac{V_N^-(N^\beta,N^t)}{V_N(N^t)}=f(\beta/t)+o(1)
       \quad(\beta/t\ge2).
\end{equation}
Here \(F,f\) are the ordinary linear-sieve functions, not the modified
upper function \(F^*\), and no distribution theorem is used in this
density calculation.

The precise classical input is the cubic Rosser positive-exit
calculation at dimension one and parameter \(2\): its compactly
convergent odd and even exit series give \(F-1\) on \(s>1\) and
\(1-f\) on \(s\ge2\), respectively. The recurrence is
\[
 a_1(s)=(3-s)_+/s,\qquad
 a_j(s)=s^{-1}\int_{\max(s,\,2+\epsilon_j)}^{2+j}a_{j-1}(u-1)\,du,
 \quad \epsilon_j={\bf1}_{j\ {\rm odd}},
\]
where an empty integral is zero.
Exact source locations are
\href{https://www.suzuyu1729.com/notes/NagoyaSeminar2022/RosserIwaniec.pdf}
{Suzuki, Sections 6--9, (9.1)--(9.2), Propositions 9.4 and 11.8,
and Lemma 14.7}; the latter is the one-sided estimate for the
actual raw Rosser density, not a different existential sieve weight.
For the normalization and fixed-density convergence one can also use
\href{https://mathweb.tifr.res.in/Documents/Publications/Lectures/tifr72.pdf}
{Motohashi, TIFR 72, (3.2.10) and Section 3.3}.
These standard results are external mathematical inputs. Neither
finite-depth convergence alone nor a numerical DDE table substitutes
for their convergence assertion.

Here is a direct reduction from those fixed local factors to the
target-dependent factors, strengthening the earlier project proof.
Let \(g_0(p)={\bf1}_{p>2}/(p-1)\), and normalize each rank-\(j\)
positive exit sum by \(V_N(y)\):
\[
 T_{j,N}(R,y)=\sum_{\substack{P=q_1\cdots q_j\ {\rm first\ exit}\\q_1<y}}
 g_N(P)\frac{V_N(q_j)}{V_N(y)}.
\]
For an exit containing no prime dividing \(N\), its summand equals
\begin{equation}\label{eq:own-lower-exit-deletion}
 g_0(P)\frac{V_0(q_j)}{V_0(y)}
 \prod_{\substack{q_j\le p<y\\p\mid N}}(1-g_0(p)).
\end{equation}
The interval is \([q_j,y)\), consistent with the strict cutoff in
\(V_N\); the factor at \(q_j\) is absent because \(q_j\nmid N\).
Other exits disappear. Every surviving multiplier lies in \([0,1]\),
and hence \(0\le T_{j,N}\le T_{j,0}\) exactly, at every rank.
The full exit identities therefore give
\[
 V_N^+/V_N\le V_0^+/V_0,\qquad V_N^-/V_N\ge V_0^-/V_0.
\]

For each fixed rank the preceding squares imply, writing
\(A_i=q_1\cdots q_i\), that \(A_i^2<RA_{i-1}\) and
\(A_i<R^{1-2^{-i}}\). Failure of the next cube forces
\(q_j^3\ge R/A_{j-1}>R^{2^{1-j}}\). Thus its prime variables all
exceed \(N^{c_j}\) uniformly on a compact parameter set.
Removing primes dividing \(N\) in this finite-dimensional integral
costs \(O(N^{-c_j}\log N)\); the multiplier in
\eqref{eq:own-lower-exit-deletion} is also \(1+o(1)\).
Prime harmonic integration gives \(T_{j,N}=a_j(s)+o(1)\).
For the upper density, retain finitely many positive exits for the
lower limit and use \(V_N^+/V_N\le V_0^+/V_0\) for the upper limit.
For the lower density reverse these two roles.
The fixed-\(g_0\) raw bound plus compact convergence of the displayed
classical series closes both limits, first \(N\to\infty\), then the
fixed depth \(j\to\infty\). There is no uncontrolled infinite-rank
limit uniform in \(N\).

There is an especially explicit check for the missing low range.
For \(1<s=\beta/t\le3\), the first upper cube forces every supported
prime below \(R^{1/3}\), and hence
\[
 V_N^+(R,y)=V_N^+(R,R^{1/3})\quad\hbox{exactly},\qquad
 \frac{V_N(R^{1/3})}{V_N(y)}=\frac3s+o(1).
\]
Using the equality at \(s=3\) gives
\(V_N^+(R,y)/V_N(y)=(3/s)F(3)+o(1)=2e^\gamma/s+o(1)=F(s)+o(1)\).
This uses the same raw coefficients, not a subtraction of upper
estimates. It covers every negative fifth-node parameter \(1<s<2\).

\paragraph{7. The signed subtree identities and the final implication.}
For a fixed retained prefix \(P_r\), residual primes are strictly less
than \(p_r\), whereas all prefix primes are at least \(p_r\).
Thus \((P_r,d)=1\) for every residual divisor and
\(g_N(P_rd)=g_N(P_r)g_N(d)\). Requiring \((P_r,N)=1\) is part of the
prime labels. There are no additional deleted prefix factors in the
residual local function.

The three replacements have the following exact density differences:
\begin{align}
 \Delta V_3&=\sum_{P_3}g_N(P_3)
  [V_N^+(N^{\nu_\rho}/P_3,p_3)-V_N^+(N^{\lambda_3}/P_3,p_3)],\nonumber\\
 \Delta V_4&=\sum_{P_4}g_N(P_4)
  [V_N^-(N^{\lambda_4}/P_4,p_4)-V_N^-(N^{\lambda_3}/P_4,p_4)],\label{eq:own-lower-subtree-telescope}\\
 \Delta V_5&=\sum_{P_5}g_N(P_5)
  [V_N^+(N^{\lambda_4}/P_5,p_5)-V_N^+(N^{\lambda_5}/P_5,p_5)].\nonumber
\end{align}
They compare successive stages of one tree, so a descendant is not
counted twice. Negative ranks three and five have upper residual
parity; positive rank four has lower residual parity. At rank three
the old history implies \(\nu_\rho-S_3>t_3\); (A) implies
\(\lambda_3-S_4>2t_4\), and
\eqref{eq:own-lower-fifth-square} implies \(\lambda_4-S_5>t_5\).
Consequently every use of \eqref{eq:own-lower-ordinary-equalities}
is in its actual domain.

The raw density of this tree is its ordinary lower start plus
\eqref{eq:own-lower-subtree-telescope}. Ordinary density equality,
prime harmonic integration, and
\[
 \pi(N)V_N(N^t)/{\cal H}_N=(2e^{-\gamma}/t)+o(1),\qquad
 {\cal H}_N=C_NN/\log^2N
\]
give its coefficient \(4\bar f(\nu_\rho/w)/w+I_3(\Pi)+I_4(\Pi)+I_5(\Pi)\),
with the three kernels and domains stated in the joint approximation
lemma and with the actual finite caps \(\Pi\).
The factor four is \(2e^{-\gamma}\cdot2e^\gamma\).
Combining the pointwise lower weight, its density lower estimate and
\eqref{eq:own-lower-signed-prime-remainder} proves, for every fixed
legal buffered finite policy and every prescribed compiler error
\(\eta>0\),
\[
 S({\cal A}_N;P(N),N^w)\ge
 \left\{\frac4w\bar f(\nu_\rho/w)+I_3(\Pi)+I_4(\Pi)+I_5(\Pi)
                 -\eta+o_N(1)\right\}{\cal H}_N.
\]
All finite choices precede \(N\). Afterwards the joint approximation
lemma removes the fixed cap/mesh/cutoff errors using its convergent
tails. Thus the exact dependency graph is classical density and
dispersion theorems, then this finite compiler, then joint finite
approximation, then the certified ideal integral. It has no arrow
back from a numerical Chen constant or from the approximation lemma
to the compiler. The two ordinary starts enter the eleven-term
functional with weights \(3/4\) and \(1/4\); a change of the L2 start
must still remove its old H integral over the changed interval.

\paragraph{8. Cross-check and scope.}
The newly supplied \texttt{own-route-uniform-singular-series-20261009.tex}
has the correct exact Euler identity: the tail of the twin-prime
constant is inverted, and the omitted large divisors of \(N\) appear
with factors \(1-1/(p-1)\). Its logarithmic tail bound
\((1+\log N/\log z)/(z-2)\) follows from the stated telescoping
integer sum and the distinct-divisor bound. It gives the normalization
above uniformly for every fixed positive exponent. It neither proves
nor is needed to assume the residual density equalities; those are
proved separately in paragraph 6.

The static final implication in
\texttt{own-route-final-transfer-20261009.tex} is consistent with this
scope. Its definition excludes the unit \(N-p=1\), counts a repeated
prime in \(\Omega\), and puts squares in the bad set before discarding
their \(O(\sqrt N)\) count in an estimate. Thus its good-count identity
is exact. Two losses of \((c_*-c)/4\) leave a strict liminf above \(c\).
Subtracting the uniform bad-count limsup and then absorbing the
remaining \(o(1)\) is valid. The ordered semiprime integral is
\(\int_{b/(1+b)}^{1/2}dt/[t(1-t)]=\log(1/b)\), with no extra factor.
This checks the logical transfer, not a fresh computation of its
saved logarithm enclosures or the table-insertion interface.

No counterexample was found for this precisely licensed finite-policy
class. The significant gap repaired here is the formerly compressed
lower compiler/density connection, particularly the actual \(s<2\)
upper residual and the target-uniform passage from classical raw
densities. This is a scoped analytic derivation using the stated
external classical and dispersion theorems. It does not independently
reprove L23's spectral estimates or the rough-SW lemma, validate all
double-sieve descendants, replay any of the 1557 frozen numerical
bindings, establish an effective threshold for \(N\), or assert that
the whole research objective or a priority claim has been verified.
